% META: {"id":"1SPE-SUITES-EX-020","chapitre":"1SPE-SUITES","type_objet":"exercice","capacites":["1SPE-SUITES-C7"],"capacites_codes":["C7"],"methodes":["M7"],"parcours":1,"competences":["calculer","raisonner"],"duree_min":10,"mode_creation":"ex_nihilo","sources_inspiration":[],"parametres_sympy":{"u0":2,"n_cible":5},"fichier_tex":"chapitres/1SPE-SUITES/exercices/1SPE-SUITES-EX-020.tex","corrige_tex":"chapitres/1SPE-SUITES/corriges/1SPE-SUITES-CO-020.tex","status":"approved","origin":"generated","status_history":["generated","audited","accepted","approved"],"release_acceptance":"RELEASE_OWNER_BATCH_ACCEPTANCE","acceptance_closure_digest":"sha256:9b3ccf9a81c5520fbb7e03b7b2d3e7bf2057a02834b6d908bf3be5f49f3b553f","acceptance_receipt_digest":"sha256:4fad86f0282e0fa4e0419cca1ad6379ae7af5a280c04a4a90880f90a1c044242"}
% BEGIN-VERIFY
% # Simulation de l'algorithme : u0=2, u_{n+1} = 2*u_n + 3
% u = 2
% for k in range(5):
%     u = 2*u + 3
% # u apres 5 iterations = u_5
% assert u == 157
% # Verification pas a pas
% vals = [2]
% for k in range(5):
%     vals.append(2*vals[-1] + 3)
% assert vals == [2, 7, 17, 37, 77, 157]
% END-VERIFY
\begin{exercice}{1SPE-SUITES-EX-020}{1}{10}
On considère la suite $(u_n)$ définie par $u_0 = 2$ et $u_{n+1} = 2u_n + 3$.

On souhaite calculer $u_5$ à l'aide d'un algorithme.

Voici le code Python correspondant :

\lstinputlisting[language=Python]{chapitres/1SPE-SUITES/python/1SPE-SUITES-EX-020.py}

\begin{enumerate}
  \item Recopier et compléter le tableau de valeurs suivant en simulant l'algorithme à la main :

  \medskip
  \begin{tabular}{|c|c|c|c|c|c|c|}
    \hline
    $n$ & $0$ & $1$ & $2$ & $3$ & $4$ & $5$ \\
    \hline
    $u_n$ & $2$ & \dots & \dots & \dots & \dots & \dots \\
    \hline
  \end{tabular}
  \medskip

  \item Quelle valeur affiche \texttt{print(u)} à la fin de l'algorithme ?
  \item Que faut-il modifier dans le code pour calculer $u_{10}$ à la place de $u_5$ ?
\end{enumerate}
\end{exercice}
